\documentclass[12pt]{letter}

\usepackage{amsmath}
\usepackage{amssymb}
\usepackage{hyperref}
\usepackage{geometry}
\usepackage{longtable}
\geometry{margin=1.25in}

\newcommand{\heading}[1]{\bigskip\noindent{\large\bfseries #1}\par\medskip}
\newcommand{\subheading}[1]{\medskip\noindent{\bfseries #1}\par\smallskip}
\newcommand{\referee}[1]{\begin{quote}\itshape #1\end{quote}}

\signature{A.\ G.\ Kim \\ P.\ E.\ Nugent \\ L.\ Wang}
\address{Physics Division\\
Lawrence Berkeley National Laboratory\\
Berkeley, California 94720, USA\\
\texttt{agkim@lbl.gov}}

\begin{document}

\begin{letter}{Editor\\
\textit{Physical Review D}\\
American Physical Society}

\opening{Dear Editor,}

We thank the Referee for a careful and constructive report on manuscript
DT18474/Kim.  We have acted on all but one of the Referee's suggestions,
which has led to significant improvements to the article; the exception is
the request for an estimate of $\alpha$ from \textit{Fermi}-LAT imaging data,
and we explain in the table below why that measurement lies outside what the
present framework can do.  The manuscript is now submitted for
reconsideration.

In this letter we summarize the changes made to the manuscript.  We submit a
version of the article compiled with the \LaTeX\ ``changes'' package to show
all the changes that were made.

The Referee's primary concern was the lack of connection to existing
literature on the subject.  We have filled two significant gaps.  We discuss
the measurement of $\alpha$ using GRB\,221009A when instrumental acceptance
in the imaging channel is taken into account, and we consider the analogue of
those effects in our spectroscopic probe.  The summary is that our results
are not inconsistent with those of Steinbring and Steinbring \& Ng, and that
the analogous effects do not affect our results.  This material is in new
paragraphs in the Introduction, in Section~IV\,B, and in the Conclusions.

We add to the Introduction a discussion of the spectral-channel analysis of
Cooke et al.  We emphasize that our distortion is derived from the
first-order coherence of the field, whereas Cooke et al.\ postulate a
convolution kernel --- a heuristic that, as those authors note, has not been
demonstrated from a model of spacetime quantization --- and that the two
predict fundamentally different effects on spectra.

Two further additions were made that the Referee did not request.  First, we
added a paragraph in Section~II\,A motivating our time-warped model for the
electric field.  Secondly, in the same section and in the Conclusions, we
discuss the difference in predicted phase jitter for light propagating along
rays versus through volumes (the analyses of Chen, Wen \& Ma and of Lee,
Zurek \& Chen), with a matching condition recorded in the Conclusions.
Thirdly, Section~III\,B now identifies the carrier--envelope term of the
closed form as the finite-linewidth correction to the per-mode core weight,
which bears on comment~4 below.
Smaller changes --- a M\"ossbauer recoil-free-fraction analogy, tightened
wording on the PKS\,1413$+$135 exclusion, and the flux-conservation
rewording described below --- are visible in the changes-marked copy.

A related change of emphasis: our effect is first order in the electric
field, entering through the first-order coherence $g^{(1)}(\tau)$.  Where the
submitted text spoke of conservation of photon number --- and in one place of
the metric exchanging energy with the photon --- the revision states instead
the property the bounds actually use: the modulation $e^{i\omega\delta t}$
has unit modulus, so the integrated flux is conserved exactly and the foam
redistributes it in frequency.

We have compiled a table of the complete set of Referee comments and our
responses to them.

\medskip

\begin{longtable}{@{}p{0.03\textwidth}p{0.42\textwidth}p{0.49\textwidth}@{}}
\hline
\textbf{\#} & \textbf{Referee comment} & \textbf{Response} \\
\hline
\endfirsthead
\hline
\textbf{\#} & \textbf{Referee comment} & \textbf{Response} \\
\hline
\endhead

1 & ``In optical photometry, this is typically refered to as a `halo.'
Mentioning this previous nomenclature adopted by others may help show how
their work compares.''
& Adopted.  Section~IV\,B now identifies the scattered fraction as the
incoherent halo of optical photometry, and names the pedestal as its
spectroscopic counterpart. \\[2ex]

2 & ``both AGNs and GRBs have to be viewed (depending on wavelength) as
potentially extended sources; that should be mentioned here.''
& Added.  Section~II\,B states that the angular extent of the source does not
enter, the sightlines superposing in intensity.  A new Appendix~A proves
this, and also treats collection over an extended aperture. \\[2ex]

3 & ``not merely scattering to the horizon \ldots\ any that don't fall within
the resolution limits of their beams, as assumed here.  The authors can more
clearly state that assumption at the outset.''
& Done.  A new paragraph in the Introduction states that the article treats
spectral distortion and does not model deflection out of the instrumental
field of view, and shows that deflection drops out of each of the three
observables used. \\[2ex]

4 & ``does not sufficiently address how real telescopes actually sense
photons \ldots\ treated as `flat' or uniform scattering outside the intrinsic
instrumental PSF assuming a single, monochromatic strength of blurring
\ldots\ It does not report an observed averaged effect over the bandpass, as
a real telescope would.''
& Addressed in Sections~IV\,B and~V\,A.  The instrumental line-spread
function is now given and shown not to enter: the BGO resolution near an MeV
is $\approx\!10\%$ FWHM, improving toward higher energies across the line's
$6$--$37$\,MeV span, so $\sigma_{\rm LSF}/E\lesssim0.04$, well below the
measured fractional width, and the quoted widths are already
response-corrected; the TeV probe counts events rather than measuring a
profile; and the FIRAS fit runs over $43$ channels, carrying the apodized
instrument line shape through its covariance.  The fraction of redistributed
flux a finite band retains is $f_{\rm band}\approx4\times10^{-17}$ for LHAASO
and $\sim\!10^{-31}$ for FIRAS.  On the monochromatic point in particular:
$e^{-\sigma_\phi^2}$ is indeed a per-mode weight, but the closed form of
Section~III\,B is not evaluated at a single frequency --- it follows from the
full line profile, and averaging the per-mode weight across that profile
reproduces exactly the $\omega_0^2\sigma_\Delta^4$ term the closed form
already carries.  The line average is built in rather than assumed
away. \\[2ex]

5 & ``that result does not preclude the \emph{range} of values which [12]
finds as consistent with observations: $0.650\leq\alpha\leq0.735$ \ldots\
those results are, in effect, a first \emph{detection}, effectively dismissed
here \ldots\ which is incorrect.''
& Accepted; no dismissal was intended.  The Introduction and Conclusions now
present that range as a two-sided bracket rather than a one-sided limit,
state that our two pedestal-independent floors lie below it and exclude no
part of it, and show that combined with it our conditional TeV bound
constrains $\tau_{\rm foam}$ from below rather than excluding
$\alpha=2/3$. \\[2ex]

6 & ``the later work of Steinbring and Ng (following the Steinbring, 2015,
ApJ, 802, 38 approach; here uncited).''
& Cited in the Introduction, at the point where the treatment of telescope
optics is discussed. \\[2ex]

7 & ``This also corrects a flaw of not accounting for more-blurred, bluer
photons reported in Steinbring, 2007, ApJ, 655, 714 (also uncited).''
& Cited, and the point is now stated.  The delay is achromatic but the phase
it produces is not, so a shorter wavelength turns the same path-length jitter
into a larger phase and the bluest photons are blurred most; accounting for
their loss beyond the instrumental acceptance is what permits a larger
$\alpha$ at fixed image sharpness. \\[2ex]

8 & ``this should properly be referred to as the Ng-Steinbring relation, I
believe, as Ng-Perlman would have been ruled out by GRB observations.''
& Partly adopted.  Equation~(1) is a variance relation due to Ng and
collaborators and is common ground; what separates the analyses built on it
is the mapping from that variance to an observable.  We have therefore
dropped the attributive label from the relation throughout, rather than
substituting a different one.  Section~IV\,B sets out the two mappings --- a
Strehl ratio, and an angular size subject to instrumental acceptance --- and
notes that this work uses neither.  We agree with the physical point behind
the request: a single Strehl-type angular dispersion at TeV energies would
scatter sources across the sky and no GRB would be localizable, so
instrumental acceptance must be modeled; the revised Section~IV\,B states
this. \\[2ex]

9 & ``Acknowledgement of that, and that the Steinbring (2023) result would
then be a first detection of foam-induced spatial blurring --- implying that
the in-orbit observed Fermi LAT PSF is, in fact, broadened by the effects of
holographic quantum foam.''
& Added to the Introduction, reported in those authors' own terms: the
bracket endpoints and their instrumental meaning ($\alpha=0.650$ scaled to
the field of view, $\alpha=0.735$ the weakest-blurring case scaled to the
nominal resolution), the holographic $\alpha=2/3$ inside the bracket, and their
finding that LAT images are not diffraction-limited but scale as
foam-induced blurring predicts. \\[2ex]

10 & ``Can the authors estimate, if incorporating the Ng-Steinbring approach
to alpha on the Fermi LAT data, that they too find a value of alpha$>$0.65 is
favoured?''
& We regret that we cannot.  Such an estimate would require adopting one of
the imaging mappings that this work does not use: the LAT records arrival
directions rather than spectra, so the analysis would be an imaging analysis
rather than a spectroscopic one.  The suppression factor $e^{-\sigma_\phi^2}$
here is derived from the first-order coherence itself and coincides with the
Strehl factor only because both are the Gaussian average $\langle
e^{i\phi}\rangle$, taken across the aperture there and across the lag here;
no angular inference enters.  What we can say, and now do say, is that our
floors lie below $0.65$ and are consistent with a larger value being
favoured. \\[2ex]

11 & ``linewidth-jitter blurring by foam has already been looked for in the
optical by Cooke et al., 2020, MNRAS, 494, 4 \ldots\ there should be some
discussion of how this current work compares to that, and it should be cited
here.''
& Cited, with a new paragraph in the Introduction comparing the two.  Their
kernel is postulated; ours is derived from the first-order coherence of the
field propagating through a fluctuating metric.  The two also differ in what
they predict: their Gaussian kernel keeps the displaced flux within the line,
whereas here it is displaced to $\sim\!1/\tau_{\rm foam}$ and leaves the
band. \\[2ex]

12 & ``the reference to [13] \ldots\ is now published.  It should be \ldots\
Proceedings of Science \ldots\ (CORFU2025) 205.''
& Updated in the bibliography. \\[2ex]

13 & ``the paper \ldots\ does \emph{not} move (although is consistent with)
the currently-known limits of the alpha-value: a holographic value of 2/3 is
still favoured.''
& Agreed, and now stated.  The spectroscopic floors do not move the imaging
limit; what they add is an independent route to the same exponent, through a
different observable and with unrelated systematics, and the two agree. \\
\hline
\end{longtable}

\medskip

We believe that the manuscript is substantially improved by these changes,
and we are grateful to the Referee for the care taken with the report.  We
hope that the revised version is now suitable for publication in
\textit{Physical Review D}, and we would of course be glad to address any
remaining concerns.

\closing{Sincerely,}

\end{letter}

\end{document}
